\subsection{What the counterexample says about L1}

The failure of the main representability criterion does not by itself
refute the separate constructibility assertion L1.  In fact its
counterexample has an open modularity locus on every permitted test.

\begin{proposition}[The modularity locus of the henselized counterexample]
\label{A:hensel-L1-open}
\status{Proved}
Let $S'$ be a locally noetherian $\F_p[t]$-scheme, let $T$ be
locally of finite type over $S'$, and let $\eta\in P(T)$ for
$P(R)=\alpha_p(R^h_{(t)})$.  Set
\[
 U_0=\{s\in S':t\text{ is nilpotent in }\mathcal O_{S',s}\},
 \qquad U_1=D(t).
\]
Both subsets are open.  No point of $T$ over
$S'\setminus(U_0\cup U_1)$ is locally modular for $P_{S'}$.
On $T_{U_0}$, respectively $T_{U_1}$, local modularity is exactly
the local-ring-isomorphism condition for the morphism defined by
$\eta$ to $\alpha_{p,U_0}$, respectively to $U_1$.
Consequently $M(T,\eta)$ is open.  It is therefore locally
constructible, and its intersection with every quasi-compact
noetherian open subscheme of $T$ is constructible.
\end{proposition}
\begin{proof}
Nilpotence of a germ of the function $t$ means that some power of
$t$ vanishes on a neighborhood, so $U_0$ is open; $U_1$ is open
by definition.  On the first open the functor is $\alpha_p$, and
on the second it is terminal.  These assertions can be checked
on neighborhoods where $t$ is nilpotent or invertible and then
glued.  They give the stated represented descriptions.

Suppose, for a contradiction, that $z\in T$ is locally modular
and that $s\in S'$ below it has $t$ nonunit and nonnilpotent in
$A_0=\mathcal O_{S',s}$.  The functor $P$ has exactly one element
over every field, so the minimal field-of-definition condition
in actual local modularity gives $\kappa(z)=\kappa(s)$.
On all local Artinian $A_0$-algebras centered at its maximal ideal,
$t$ is nilpotent and $P=\alpha_p$.  Thus the universal property
at $z$, restricted to these Artinian tests, identifies complete
local representing rings:
\[
 \widehat{\mathcal O_{T,z}}
       \simeq \widehat A_0[x]/(x^p).
\]
In particular the local ring of the fiber at $z$ has dimension
zero: its completion is $\kappa(s)[x]/(x^p)$.
Since $T\to S'$ is locally of finite type, it is quasi-finite
at $z$.

Put $A=\widehat A_0$.  This is a complete noetherian local ring,
and $t$ remains nonnilpotent because $A_0\to A$ is faithfully
flat.  Let $z'$ be the point of $T\times_{S'}\Spec A$ over
$z$ and the closed point of $\Spec A$; it is unique because the
two residue fields just identified are equal.
The pulled-back pair is actually locally modular at $z'$ over
$A$.  Indeed every local $A$-test centered there is a local
$S'$-test centered at $z$; the unique map to $T$ supplied by
actual modularity, together with its prescribed map to $\Spec A$,
is the required unique map to the fiber product.  The residue
field condition is unchanged.

Choose an affine quasi-finite neighborhood of $z'$ over $A$.
It is separated, and Zariski's Main Theorem embeds it as an open
subscheme of a finite $A$-scheme.  Since $A$ is henselian local,
the latter is a finite disjoint union of spectra of finite local
$A$-algebras.  The factor through $z'$ is entirely contained in
the chosen open: an open subset of a local scheme containing its
closed point is the whole scheme.  We have therefore produced
a finite local $A$-algebra with an actual locally modular
element.  Proposition~\ref{A:hensel-remaining-gate} proves that no
such pair exists when $t$ is nonunit and nonnilpotent in a complete
noetherian local base.  This is the desired contradiction.

It remains to check openness on $T_{U_0}$ and $T_{U_1}$.
For either represented restriction, actual local modularity means
that the corresponding local-ring homomorphism is an isomorphism.
The source and target are locally of finite type over the
locally noetherian base, so their comparison morphism is locally
of finite presentation.  Its local-ring-isomorphism locus is
open: at such a point the morphism is \emph{\'etale}; the inverse
map on local schemes spreads by finite presentation to a section
over a neighborhood in the target, and that section has open
image after restricting to the \emph{\'etale} locus.
This is the same local argument as in
Proposition~\ref{A:8-literal}.
The two resulting open subsets of $T$ are all of $M(T,\eta)$
by the preceding exclusion.  This proves the proposition.
\end{proof}

This proposition decides the logical relation to the new main
counterexample: it satisfies the conclusion asked for in L1.
The general assertion requires a separate argument, given in
Theorem~\ref{A:L1-general} below; it does not follow merely
from this counterexample-specific computation.
